% scon_en_ebnf.tex
% SCON 1.0 — Specification, Benchmarks, and Token Efficiency Analysis
% arXiv submission (pdflatex compatible)

\documentclass[11pt, a4paper]{article}

\usepackage[a4paper, top=2.5cm, bottom=2.5cm, left=2cm, right=2cm]{geometry}
\usepackage[T1]{fontenc}
\usepackage[utf8]{inputenc}
\usepackage{hyperref}
\usepackage{amsmath}
\usepackage{amssymb}
\usepackage{booktabs}
\usepackage{xcolor}
\usepackage{listings}
\usepackage{syntax}

% EBNF listing style
\lstdefinelanguage{EBNF}{
  morekeywords={},
  morecomment=[l]{(*},
  morestring=[b]",
  literate=
    {::=}{{$::=$}}3
    {|}{{$\mid$}}1
    {[}{{$[$}}1
    {]}{{$]$}}1
    {\{}{{$\{$}}1
    {\}}{{$\}$}}1
    {(}{{$($}}1
    {)}{{$)$}}1
    {*}{{$*$}}1
    {+}{{$^+$}}1
}

\lstdefinestyle{ebnf}{
  language=EBNF,
  basicstyle=\ttfamily\small,
  keywordstyle=\bfseries,
  commentstyle=\itshape\color{gray},
  stringstyle=\color{teal},
  frame=none,
  breaklines=true,
  columns=fullflexible,
  keepspaces=true,
  xleftmargin=1em,
  aboveskip=0.5em,
  belowskip=0.5em,
}

\lstdefinestyle{scon}{
  basicstyle=\ttfamily\small,
  frame=single,
  backgroundcolor=\color{gray!8},
  breaklines=true,
  columns=fullflexible,
  keepspaces=true,
  xleftmargin=1em,
  aboveskip=0.5em,
  belowskip=0.5em,
}

\title{\textbf{SCON: A Textual Serialization Format with Structural Deduplication,}\\
\textbf{Tabular Encoding, and Token-Efficient Representation}}
\author{Gustavo Moura \\
\small CIFRID \\
\small \texttt{moura@cifrid.com}}
\date{}

\begin{document}
\maketitle

\begin{abstract}
SCON (Structured Compact Object Notation) is a human-readable,
text-based serialization format designed as a compact alternative
to JSON for payload-sensitive workloads.
Tabular encoding eliminates repeated keys in uniform arrays,
and optional structural deduplication (via Merkle-style hashing)
factorizes recurring subtrees into named references.
Benchmarks on canonical datasets show payloads 13--29\% smaller than
JSON without deduplication (up to 66\% with deduplication), while
reducing transport-layer compression CPU by 8--53\%.
A break-even analysis demonstrates that byte savings offset the
additional parsing cost on links below $\sim$100\,Mbps, a regime
common in IoT telemetry, mobile networks, and satellite backhaul.
On an industrial telemetry workload (10K sensor readings), SCON's
tabular mode reduces payload size by 53\% vs JSON.
Structural deduplication also reduces LLM tokenization cost:
on an OpenAPI specification, SCON produces 64\% fewer tokens than
the equivalent JSON under the \texttt{cl100k\_base} tokenizer,
enabling more efficient use of LLM context windows.
This paper presents the formal grammar (EBNF), reference
implementations in Rust, PHP, and JavaScript, and a quantitative
evaluation of size, speed, and token efficiency trade-offs.
\end{abstract}

\medskip
\noindent\textbf{Keywords:}
serialization format,
JSON alternative,
structural deduplication,
tabular encoding,
payload optimization,
IoT serialization,
industrial telemetry,
SCADA data compression,
sensor data format,
LLM context optimization,
prompt token reduction,
dataset serialization,
edge computing,
bandwidth-constrained networks

\tableofcontents
\newpage

%% ============================================================
\section{Introduction}
\label{sec:intro}

JSON~\cite{rfc8259} is the dominant data interchange format for web APIs,
configuration files, and inter-process communication. Its simplicity and
ubiquitous tooling support are well established. However, JSON exhibits
structural inefficiencies when serializing repetitive data:

\begin{itemize}
\item \textbf{Key repetition.} In arrays of homogeneous objects, every
      element repeats the full key set. A 10\,000-row sensor dataset
      repeats each field name 10\,000 times.
\item \textbf{No structural reuse.} Identical subtrees (e.g., shared
      response schemas in an OpenAPI specification) must be serialized
      in full at every occurrence.
\item \textbf{Readability--size trade-off.} JSON requires pretty-printing
      (2--4$\times$ size increase) to be human-readable; the minified
      form is compact but opaque.
\end{itemize}

\noindent SCON (Structured Compact Object Notation) addresses these
limitations with three mechanisms:

\begin{enumerate}
\item \textbf{Tabular encoding:} Arrays of uniform objects are serialized
      with a single header row followed by value-only data rows,
      eliminating $R \times K$ key repetitions.
\item \textbf{Structural deduplication:} A Merkle-style fingerprinting
      pass identifies repeated subtrees and replaces them with named
      schema references (\texttt{@s:name}), reducing both bytes and
      LLM tokens.
\item \textbf{Indentation-based structure:} The standard form is
      human-readable without pretty-printing overhead; a reversible
      minification (semicolon-based depth encoding) provides a compact
      wire format.
\end{enumerate}

\subsection{Contributions}

This paper makes the following contributions:

\begin{enumerate}
\item A formal grammar specification (EBNF) for SCON~1.0, covering
      standard, minified, and schema-augmented representations
      (Sections~\ref{sec:notation}--\ref{sec:summary}).
\item Reference implementations in three languages (Rust, PHP,
      JavaScript) with cross-language roundtrip verification on
      canonical fixtures.
\item A quantitative evaluation of payload size, encode/decode
      throughput, transport-layer compression cost, and LLM token
      efficiency against JSON
      (Sections~\ref{sec:performance}--\ref{sec:roadmap}).
\item A break-even analysis identifying the network bandwidth threshold
      at which SCON's byte savings offset its parsing overhead.
\end{enumerate}

\subsection{Scope and Non-Goals}

SCON is designed for \textbf{structured data serialization} in
network-bound and storage-bound scenarios: API payloads, configuration,
telemetry, and LLM context packing. It is \emph{not} a replacement for:

\begin{itemize}
\item Binary formats (CBOR, MessagePack, Protocol Buffers) where
      human readability is not required.
\item General-purpose programming language data models (JSON Schema
      validation, static type systems).
\item Document markup languages (XML, HTML).
\end{itemize}

%% ============================================================
\section{Related Work}
\label{sec:related}

Table~\ref{tab:related} compares SCON against established serialization
formats across twelve features relevant to payload optimization and
developer ergonomics.

\begin{table}[ht]
\centering
\caption{Feature comparison across serialization formats.
         $\bullet$ = yes, $\circ$ = partial, --- = no.}
\label{tab:related}
\small
\begin{tabular}{l cccccccc c}
\toprule
Feature & JSON & JSON5 & HJSON & YAML & TOML & Ion & CBOR & MsgPk & \textbf{SCON} \\
\midrule
Human-readable       & $\bullet$ & $\bullet$ & $\bullet$ & $\bullet$ & $\bullet$ & $\bullet$ & --- & --- & $\bullet$ \\
Comments             & --- & $\bullet$ & $\bullet$ & $\bullet$ & $\bullet$ & $\bullet$ & --- & --- & $\bullet$ \\
Built-in schemas     & --- & --- & --- & --- & --- & $\bullet$ & $\circ$ & --- & $\bullet$ \\
Structural dedup     & --- & --- & --- & $\circ$ & --- & --- & --- & --- & $\bullet$ \\
Tabular encoding     & --- & --- & --- & --- & --- & --- & --- & --- & $\bullet$ \\
Import mechanism     & --- & --- & --- & --- & --- & --- & --- & --- & $\bullet$ \\
Minified form        & $\bullet$ & $\bullet$ & --- & --- & $\circ$ & $\bullet$ & $\bullet$ & $\bullet$ & $\bullet$ \\
Unquoted keys        & --- & $\bullet$ & $\bullet$ & $\bullet$ & $\bullet$ & $\bullet$ & --- & --- & $\bullet$ \\
Canonical form       & $\circ$ & --- & --- & $\circ$ & --- & $\bullet$ & $\bullet$ & $\circ$ & $\bullet$ \\
Binary variant       & --- & --- & --- & --- & --- & $\bullet$ & $\bullet$ & $\bullet$ & --- \\
\bottomrule
\end{tabular}
\end{table}

\noindent \textbf{Key differentiators:}

\begin{description}
\item[vs.\ YAML anchors.]
  YAML supports manual node reuse via anchors (\texttt{\&}/\texttt{*}),
  but the author must identify and annotate duplicates by hand.
  SCON's \texttt{autoExtract} detects structural repetition
  algorithmically via Merkle fingerprinting, requiring zero annotation.

\item[vs.\ binary formats (CBOR, MessagePack, Ion).]
  Binary formats achieve smaller payloads and faster parsing at the
  cost of human readability. SCON occupies a different point in the
  design space: text-based and human-readable, yet 13--66\% smaller
  than JSON through structural techniques rather than binary encoding.

\item[vs.\ JSON5/HJSON/TOML.]
  These formats improve JSON ergonomics (comments, trailing commas,
  unquoted keys) but do not address payload size. None offers
  tabular encoding, structural deduplication, or a minified
  wire format.

\item[vs.\ Amazon Ion.]
  Ion provides both text and binary representations with a schema
  language (ISL). SCON's schema definitions are embedded in the
  document itself and support reference patching (overrides),
  enabling structural composition absent in ISL.
\end{description}

%% ============================================================
%% FORMAL GRAMMAR begins here
%% ============================================================

\section{Formal Grammar --- Notation Conventions}
\label{sec:notation}

\begin{tabular}{ll}
\toprule
Symbol & Meaning \\
\midrule
\texttt{::=}   & Production rule definition \\
\texttt{|}     & Alternation (ordered choice) \\
\texttt{[ ]}   & Optional (zero or one) \\
\texttt{\{ \}} & Repetition (zero or more) \\
\texttt{( )}   & Grouping \\
\texttt{" "}   & Terminal string literal \\
\texttt{? ?}   & Special sequence (prose description) \\
\texttt{-}     & Exception (set difference) \\
\bottomrule
\end{tabular}

\medskip
\noindent Whitespace within productions is not significant unless enclosed in terminals.
The notation \textsc{INDENT}($n$) denotes exactly $n$ space characters (\texttt{U+0020})
at the start of a line. Tabs are not permitted for indentation.

%% ============================================================
\section{Document Structure}
\label{sec:document}

\begin{lstlisting}[style=ebnf]
document       ::= [ header NL ] { directive NL } body EOF

header         ::= "#!scon/" version
version        ::= DIGIT+ "." DIGIT+

body           ::= value
\end{lstlisting}

\noindent The header \texttt{\#!scon/1.0} is optional. When present, it identifies
the format version and enables tooling auto-detection. The encoder emits the header
only when explicitly requested (\texttt{includeHeader: true}).

%% ============================================================
\section{Directives}
\label{sec:directives}

\begin{lstlisting}[style=ebnf]
directive      ::= "@@" directive_name [ "(" directive_arg ")" ]
directive_name ::= ALPHA { ALPHA | DIGIT | "_" }
directive_arg  ::= ? any characters except ")" ?
\end{lstlisting}

\noindent Directives modify parser behavior for the enclosing scope:

\begin{lstlisting}[style=scon]
@@strict
@@enforce(openapi:3.1)
\end{lstlisting}

\noindent \texttt{@@strict} enables contract validation: objects instantiated from
a schema reference reject undeclared fields and require mandatory fields
(marked with \texttt{!}).

%% ============================================================
\section{Imports}
\label{sec:imports}

\begin{lstlisting}[style=ebnf]
import_stmt    ::= "@use" SP path
path           ::= ? relative or absolute file path ?
\end{lstlisting}

\noindent Imports load external schema definitions into the current
\texttt{SchemaRegistry} before body parsing begins.

%% ============================================================
\section{Values}
\label{sec:values}

\begin{lstlisting}[style=ebnf]
value          ::= primitive
               |   object
               |   array
               |   schema_ref

primitive      ::= null | boolean | number | string

null           ::= "null"
boolean        ::= "true" | "false"
number         ::= integer | float
integer        ::= [ "-" ] DIGIT+
float          ::= [ "-" ] DIGIT+ "." DIGIT+ [ exponent ]
exponent       ::= ( "e" | "E" ) [ "+" | "-" ] DIGIT+
\end{lstlisting}

%% ============================================================
\section{Strings}
\label{sec:strings}

\begin{lstlisting}[style=ebnf]
string         ::= quoted_string | unquoted_string

quoted_string  ::= '"' { char } '"'
char           ::= unescaped | escape_seq
unescaped      ::= ? any Unicode except '"', '\', control ?
escape_seq     ::= "\" ( '"' | "\" | "/" | "n" | "r"
                        | "t" | ";" | "b" | "f"
                        | "u" HEX HEX HEX HEX )

unquoted_string ::= ? non-empty sequence of characters that:
                       - does not start with '"'
                       - does not start with a reserved prefix
                       - does not contain unescaped ','
                         in tabular/inline context
                       - is trimmed of leading/trailing whitespace ?
\end{lstlisting}

\subsection{Quoting Rules}
\label{sec:quoting}

A value \textbf{must} be quoted when it:
\begin{itemize}
\item Equals a reserved literal (\texttt{null}, \texttt{true}, \texttt{false})
      and the intended type is string.
\item Contains a leading or trailing space that is semantically significant.
\item Contains the delimiter character (\texttt{,} in inline/tabular contexts).
\item Contains a semicolon (\texttt{;}) in minified representation.
\item Parses as a number but the intended type is string (e.g., \texttt{"007"}).
\end{itemize}

\subsection{Escape Sequences}
\label{sec:escapes}

\begin{tabular}{lll}
\toprule
Sequence & Character & Unicode \\
\midrule
\texttt{\textbackslash"} & Double quote & U+0022 \\
\texttt{\textbackslash\textbackslash} & Backslash & U+005C \\
\texttt{\textbackslash n} & Line feed & U+000A \\
\texttt{\textbackslash r} & Carriage return & U+000D \\
\texttt{\textbackslash t} & Horizontal tab & U+0009 \\
\texttt{\textbackslash ;} & Literal semicolon & U+003B \\
\texttt{\textbackslash b} & Backspace & U+0008 \\
\texttt{\textbackslash f} & Form feed & U+000C \\
\texttt{\textbackslash /} & Solidus & U+002F \\
\texttt{\textbackslash uXXXX} & Unicode code point & U+XXXX \\
\bottomrule
\end{tabular}

\medskip
\noindent The \texttt{\textbackslash ;} escape is unique to SCON and is required
in minified mode to distinguish a literal semicolon from the structural dedent
delimiter.

%% ============================================================
\section{Keys}
\label{sec:keys}

\begin{lstlisting}[style=ebnf]
key            ::= unquoted_key | quoted_key
unquoted_key   ::= KEY_CHAR { KEY_CHAR }
quoted_key     ::= '"' { char } '"'
KEY_CHAR       ::= ALPHA | DIGIT | "_" | "-" | "." | "$"
                 | "/" | "~" | "@"
\end{lstlisting}

\noindent Keys \textbf{must} be quoted when they:
\begin{itemize}
\item Contain spaces, colons, or brackets.
\item Match a reserved literal (\texttt{null}, \texttt{true}, \texttt{false}).
\item Start with \texttt{-} (conflicts with array item prefix).
\item Contain characters outside the \texttt{KEY\_CHAR} set.
\end{itemize}

\noindent A key followed by \texttt{!} denotes a mandatory field in schema definitions:

\begin{lstlisting}[style=ebnf]
mandatory_key  ::= key "!"
\end{lstlisting}

%% ============================================================
\section{Objects}
\label{sec:objects}

\begin{lstlisting}[style=ebnf]
object         ::= { INDENT(n) key_value NL }

key_value      ::= key ":" SP value_or_block

value_or_block ::= inline_value
               |   NL nested_block

inline_value   ::= primitive | inline_array | schema_ref

nested_block   ::= object | array
\end{lstlisting}

\noindent Objects are defined by indentation level. All key-value pairs at the same
indentation depth belong to the same object. A key followed by \texttt{:} with no
inline value introduces a nested block at depth $n+1$.

\begin{lstlisting}[style=scon]
server:
 host: localhost
 port: 8080
 tls:
  enabled: true
  cert: /etc/ssl/cert.pem
\end{lstlisting}

%% ============================================================
\section{Arrays}
\label{sec:arrays}

SCON supports three array representations, each optimized for a different data shape.

\subsection{Expanded Arrays (List Items)}
\label{sec:expanded-arrays}

\begin{lstlisting}[style=ebnf]
expanded_array ::= { INDENT(n) "- " array_item NL }

array_item     ::= primitive
               |   schema_ref
               |   inline_object_start NL expanded_object
               |   NL nested_block

inline_object_start ::= key ":" SP inline_value
\end{lstlisting}

\noindent Each item is prefixed with \texttt{-~} (dash + space). Objects within
list items begin their first key-value pair on the same line as the dash:

\begin{lstlisting}[style=scon]
users:
 - name: Alice
   role: admin
 - name: Bob
   role: viewer
\end{lstlisting}

\subsection{Inline Arrays}
\label{sec:inline-arrays}

\begin{lstlisting}[style=ebnf]
inline_array   ::= "[" count "]" ":" SP value_list
value_list     ::= inline_value { "," SP inline_value }
count          ::= DIGIT+
\end{lstlisting}

\noindent The count $N$ declares the number of elements. Values are
comma-separated on a single line:

\begin{lstlisting}[style=scon]
tags[3]: web, api, rest
ports[4]: 80, 443, 8080, 8443
\end{lstlisting}

\subsection{Tabular Arrays}
\label{sec:tabular-arrays}

\begin{lstlisting}[style=ebnf]
tabular_array  ::= tabular_header NL { INDENT(n+1) data_row NL }

tabular_header ::= key "[" row_count "]" "{" field_list "}" ":"

field_list     ::= field_name { "," SP field_name }
field_name     ::= key

row_count      ::= DIGIT+

data_row       ::= inline_value { "," SP inline_value }
\end{lstlisting}

\noindent Tabular encoding compacts arrays of homogeneous objects. The header
declares field names once; each subsequent line contains only the values
in corresponding order:

\begin{lstlisting}[style=scon]
users[3]{id, name, role}:
 1, Alice, admin
 2, Bob, editor
 3, Carol, viewer
\end{lstlisting}

\noindent \textbf{Detection heuristic.} The encoder activates tabular mode when:
\begin{enumerate}
\item All items in the array are objects.
\item All objects share the same key set (uniform structure).
\item All values are primitives (no nested objects or arrays).
\item The array contains $\geq 2$ items.
\end{enumerate}

\noindent When any condition fails---e.g., objects have optional fields,
mixed schemas, or nested values---the encoder falls back to
\emph{expanded} format (one \texttt{- key: value} block per item),
which preserves correctness but does not eliminate key repetition.
Tabular savings are therefore data-dependent: arrays of uniform,
flat objects (database rows, sensor readings, log entries) benefit
most, while sparse or polymorphic arrays see no tabular advantage.

%% ============================================================
\section{Schema Definitions}
\label{sec:schema-defs}

\begin{lstlisting}[style=ebnf]
schema_block   ::= schema_header NL INDENT(n+1) object

schema_header  ::= block_prefix ":" key ":"
block_prefix   ::= "s" | "r" | "sec"
\end{lstlisting}

\noindent Three block types share the same syntax:

\begin{tabular}{lll}
\toprule
Prefix & Name & Purpose \\
\midrule
\texttt{s:} & Schema & Reusable object template \\
\texttt{r:} & Response & Response group (API contexts) \\
\texttt{sec:} & Security & Security scheme definition \\
\bottomrule
\end{tabular}

\begin{lstlisting}[style=scon]
s:UserBase:
 id!: 0
 name!: ""
 email: null
 role: viewer
\end{lstlisting}

\noindent Fields marked with \texttt{!} are mandatory. When \texttt{@@strict} is
active, instances must include all mandatory fields.

%% ============================================================
\section{Schema References and Overrides}
\label{sec:schema-refs}

\begin{lstlisting}[style=ebnf]
schema_ref     ::= "@s:" schema_name [ SP override_block ]
               |   polymorphic_ref

schema_name    ::= ALPHA { ALPHA | DIGIT | "_" | "-" }

polymorphic_ref ::= schema_ref SP "|" SP schema_ref
                    { SP "|" SP schema_ref }
\end{lstlisting}

\subsection{Override Syntax}
\label{sec:overrides}

\begin{lstlisting}[style=ebnf]
override_block ::= { INDENT(n+1) override_entry NL }

override_entry ::= field_override | removal

field_override ::= dot_path ":" SP value_or_block
dot_path       ::= key { "." key }

removal        ::= "-" key
\end{lstlisting}

\noindent Overrides enable structural patching of schema instances:

\begin{lstlisting}[style=scon]
admin: @s:UserBase
 role: admin
 permissions:
  - manage_users
  - manage_system
 -email
\end{lstlisting}

\noindent In this example:
\begin{itemize}
\item \texttt{role: admin} overrides the default \texttt{viewer}.
\item \texttt{permissions:} adds a new nested field.
\item \texttt{-email} removes the \texttt{email} field from the instance.
\end{itemize}

\subsection{Deep Override (Dot-Path)}

\begin{lstlisting}[style=scon]
service: @s:ServiceConfig
 database.pool.max: 50
 database.pool.timeout: 5000
 -database.pool.debug
\end{lstlisting}

\noindent Dot-path notation traverses nested objects without redefining the
full subtree.

\subsection{Polymorphic Composition}

\begin{lstlisting}[style=scon]
endpoint: @s:BaseEndpoint | @s:AuthMixin | @s:CacheMixin
 path: /api/users
\end{lstlisting}

\noindent The pipe operator \texttt{|} merges multiple schemas left-to-right.
Later schemas override earlier ones on key conflicts.

%% ============================================================
\section{Minified Representation}
\label{sec:minification}

\begin{lstlisting}[style=ebnf]
minified_doc   ::= { token }

token          ::= key_value_min | array_item_min
               |   dedent_seq | quoted_string
               |   header | directive

key_value_min  ::= key ":" SP inline_value ";"

dedent_seq     ::= ";" { ";" }
\end{lstlisting}

\noindent The minification algorithm replaces indentation and newlines with a
unary semicolon encoding:

\begin{enumerate}
\item A single \texttt{;} separates sibling entries at the same depth
      (equivalent to a newline).
\item A sequence of $N$ semicolons ($N \geq 2$) represents a dedent of
      $N-1$ levels in the tree hierarchy.
\end{enumerate}

\noindent \textbf{Examples:}

\begin{tabular}{ll}
\toprule
Semicolons & Meaning \\
\midrule
\texttt{;} ($N{=}1$)  & Next sibling (same depth) \\
\texttt{;;} ($N{=}2$) & Close current level, next sibling of parent \\
\texttt{;;;} ($N{=}3$) & Close 2 levels, next sibling of grandparent \\
\bottomrule
\end{tabular}

\medskip
\noindent \textbf{Standard form:}
\begin{lstlisting}[style=scon]
server:
 host: localhost
 port: 8080
 tls:
  enabled: true
  cert: /etc/ssl/cert.pem
database:
 host: db.local
\end{lstlisting}

\noindent \textbf{Minified form:}
\begin{lstlisting}[style=scon]
server:host: localhost;port: 8080;tls:enabled: true;cert: /etc/ssl/cert.pem;;;database:host: db.local
\end{lstlisting}

\noindent The three semicolons \texttt{;;;} after \texttt{cert} dedent
2~levels: from depth~2 (\texttt{tls} children) back to depth~0
(\texttt{database} is a sibling of \texttt{server}).

\subsection{Reversibility Property}

The minification is fully reversible:

$$\texttt{expand}(\texttt{minify}(S)) \equiv S$$

\noindent because the unary encoding preserves exact tree depth at every
position. The \texttt{expand} function reconstructs indentation by
maintaining a depth counter and emitting $d$ spaces per line.

%% ============================================================
\section{Delimiter Variants}
\label{sec:delimiters}

\begin{lstlisting}[style=ebnf]
delimiter      ::= "," | "|" | TAB
\end{lstlisting}

\noindent In inline and tabular arrays, SCON supports three delimiter
characters. The default is comma (\texttt{,}). Pipe (\texttt{|}) and
tab (\texttt{TAB}) are recognized during decoding to support legacy
and TSV-compatible data sources.

%% ============================================================
\section{Comments}
\label{sec:comments}

\begin{lstlisting}[style=ebnf]
comment        ::= "#" { ? any character except NL ? } NL
\end{lstlisting}

\noindent Comments begin with \texttt{\#} and extend to the end of the line.
They are permitted at any indentation level and are discarded during parsing.
The header line \texttt{\#!scon/1.0} is a special case: the \texttt{\#!}
prefix distinguishes it from regular comments.

%% ============================================================
\section{Lexical Elements}
\label{sec:lexical}

\begin{lstlisting}[style=ebnf]
NL             ::= U+000A | U+000D U+000A
SP             ::= U+0020
TAB            ::= U+0009
INDENT(n)      ::= SP * n
DIGIT          ::= "0" | "1" | ... | "9"
HEX            ::= DIGIT | "a" | ... | "f" | "A" | ... | "F"
ALPHA          ::= "a" | ... | "z" | "A" | ... | "Z"
EOF            ::= ? end of input ?
\end{lstlisting}

\noindent \textbf{Indentation rules:}
\begin{itemize}
\item Only space characters (\texttt{U+0020}) are valid for indentation.
\item Tab characters are rejected in indentation context.
\item Each indentation level corresponds to exactly 1~space by convention,
      though the parser accepts any consistent increment.
\item The root level is depth~0 (no leading spaces).
\end{itemize}

%% ============================================================
\section{Auto-Extracted Schemas}
\label{sec:autoextract}

The \texttt{autoExtract} pipeline is not part of the grammar per se, but
produces valid SCON output by:

\begin{enumerate}
\item Traversing the data tree and computing xxHash128 fingerprints for
      each subtree.
\item Identifying structures that appear $\geq 2$ times with identical
      fingerprints.
\item Extracting repeated structures as named schema definitions
      (\texttt{s:auto\_$N$}).
\item Replacing occurrences with schema references (\texttt{@s:auto\_$N$}).
\item Pruning orphan schemas (schemas with $< 2$ active references).
\end{enumerate}

\noindent The resulting document is syntactically valid SCON with schema
blocks prepended to the body.

%% ============================================================
\section{Complete Grammar Summary}
\label{sec:summary}

For reference, the complete grammar is presented below without commentary.

\begin{lstlisting}[style=ebnf]
(* SCON 1.0 -- Complete EBNF *)

document       ::= [ header NL ] { directive NL }
                   { import_stmt NL } { schema_block NL }
                   body EOF

header         ::= "#!scon/" version
version        ::= DIGIT+ "." DIGIT+

directive      ::= "@@" directive_name [ "(" directive_arg ")" ]
import_stmt    ::= "@use" SP path

body           ::= value

value          ::= primitive | object | array | schema_ref

primitive      ::= null | boolean | number | string
null           ::= "null"
boolean        ::= "true" | "false"
integer        ::= [ "-" ] DIGIT+
float          ::= [ "-" ] DIGIT+ "." DIGIT+ [ exponent ]
exponent       ::= ( "e" | "E" ) [ "+" | "-" ] DIGIT+
number         ::= integer | float

string         ::= quoted_string | unquoted_string
quoted_string  ::= '"' { char } '"'
char           ::= unescaped | escape_seq
escape_seq     ::= "\" ( '"' | "\" | "/" | "n" | "r"
                        | "t" | ";" | "b" | "f"
                        | "u" HEX HEX HEX HEX )

key            ::= unquoted_key | quoted_key
unquoted_key   ::= KEY_CHAR { KEY_CHAR }
quoted_key     ::= '"' { char } '"'
mandatory_key  ::= key "!"

object         ::= { INDENT(n) key_value NL }
key_value      ::= key ":" SP value_or_block
value_or_block ::= inline_value | NL nested_block

expanded_array ::= { INDENT(n) "- " array_item NL }
inline_array   ::= "[" count "]" ":" SP value_list
tabular_array  ::= tabular_header NL
                    { INDENT(n+1) data_row NL }
tabular_header ::= key "[" row_count "]"
                    "{" field_list "}" ":"
array          ::= expanded_array | inline_array
                 | tabular_array

schema_block   ::= block_prefix ":" key ":" NL
                    INDENT(n+1) object
block_prefix   ::= "s" | "r" | "sec"

schema_ref     ::= "@s:" schema_name
                    [ SP override_block ]
               |   polymorphic_ref
polymorphic_ref ::= schema_ref SP "|" SP schema_ref
                    { SP "|" SP schema_ref }
override_entry ::= dot_path ":" SP value_or_block
               |   "-" key
dot_path       ::= key { "." key }

comment        ::= "#" { ? any char except NL ? } NL

(* Minified representation *)
dedent_seq     ::= ";" { ";" }
\end{lstlisting}

%% ============================================================
\section{Conformance Levels}
\label{sec:conformance}

A SCON implementation may declare conformance at one of three levels:

\begin{description}
\item[Level 1 --- Core.]
  Supports primitives, objects, expanded arrays, inline arrays, and
  minification/expansion. Sufficient for basic serialization.

\item[Level 2 --- Tabular.]
  Level~1 plus tabular array detection (encoder) and parsing (decoder).
  Required for payload-size claims.

\item[Level 3 --- Full.]
  Level~2 plus schema definitions, schema references with overrides,
  polymorphic composition, directives, imports, and \texttt{autoExtract}
  structural deduplication.
\end{description}

\noindent The Rust implementation is Level~2. The PHP and JavaScript
implementations are Level~3.

%% ============================================================
\section{Performance Trade-offs}
\label{sec:performance}

SCON's design trades parsing speed for payload compactness and human readability.
This section quantifies the trade-off using Rust-native benchmarks
(\texttt{serde\_json} vs \texttt{scon-core}), which provide the fairest
format-versus-format comparison since both sides are compiled code.

\subsection{Observed Ratios}
\label{sec:ratios}

All measurements use three canonical fixtures (OpenAPI Specs 49\,KB, Config
Records 73\,KB, DB Exports 20\,KB) at 200 iterations on the same host.
Both \texttt{serde\_json} and \texttt{scon-core} are compiled Rust;
this isolates format-specific overhead from language runtime differences.

\begin{table}[ht]
\centering
\caption{Rust-native encode/decode overhead (median, ms). Current
         implementation with two allocation patches applied
         (see Section~\ref{sec:impl-overhead}).}
\label{tab:rust-overhead}
\begin{tabular}{l rr c rr c}
\toprule
Dataset & \multicolumn{2}{c}{Encode (ms)} & Ratio
        & \multicolumn{2}{c}{Decode (ms)} & Ratio \\
        & JSON & SCON & & JSON & SCON & \\
\midrule
OpenAPI Specs   & 0.061 & 0.099 & 1.6$\times$ & 0.343 & 0.953 & 2.8$\times$ \\
Config Records  & 0.087 & 0.187 & 2.2$\times$ & 0.457 & 1.130 & 2.5$\times$ \\
DB Exports      & 0.022 & 0.075 & 3.4$\times$ & 0.091 & 0.312 & 3.4$\times$ \\
\bottomrule
\end{tabular}
\end{table}

\begin{table}[ht]
\centering
\caption{Payload size comparison (bytes).}
\label{tab:payload-size}
\begin{tabular}{l rrr r}
\toprule
Dataset & JSON & SCON(min) & SCON(dedup+min) & Saving vs JSON \\
\midrule
OpenAPI Specs   & 49\,067 & 42\,768 & 16\,600 & $-13\%$ / $-66\%$ \\
Config Records  & 73\,183 & 67\,049 & --- & $-8\%$ \\
DB Exports      & 19\,688 & 13\,917 & 13\,917 & $-29\%$ \\
\bottomrule
\end{tabular}
\end{table}

\noindent In the current implementation, SCON encoding is 1.6--3.4$\times$ slower and
decoding is 2.5--3.4$\times$ slower than \texttt{serde\_json}. Payload size is
8--29\% smaller without deduplication; up to 66\% smaller with structural
deduplication enabled. Section~\ref{sec:impl-overhead} analyses the sources
of this overhead and demonstrates that it is dominated by implementation
choices rather than format complexity.

\subsection{Break-Even Inequality}
\label{sec:breakeven}

SCON compensates its CPU overhead when the byte savings reduce transmission
time by more than the additional processing cost. For a single
request--response cycle, the net latency change is:

\begin{equation}
\label{eq:breakeven}
\Delta t = \underbrace{(\Delta t_{\text{enc}} + \Delta t_{\text{dec}})}_{\text{CPU overhead}}
         - \underbrace{\frac{\Delta B}{\mathit{BW}}}_{\text{transmission savings}}
\end{equation}

\noindent where:
\begin{itemize}
\item $\Delta t_{\text{enc}} = t_{\text{enc}}^{\text{SCON}} - t_{\text{enc}}^{\text{JSON}}$
      \quad (additional encode time),
\item $\Delta t_{\text{dec}} = t_{\text{dec}}^{\text{SCON}} - t_{\text{dec}}^{\text{JSON}}$
      \quad (additional decode time),
\item $\Delta B = B_{\text{JSON}} - B_{\text{SCON}}$
      \quad (bytes saved),
\item $\mathit{BW}$ is the effective link bandwidth.
\end{itemize}

\noindent SCON yields a net latency improvement when $\Delta t < 0$, i.e.:

\begin{equation}
\label{eq:breakeven-condition}
\boxed{\;\Delta t_{\text{enc}} + \Delta t_{\text{dec}} \;<\; \frac{\Delta B}{\mathit{BW}}\;}
\end{equation}

\subsection{End-to-End Latency by Network Scenario}
\label{sec:e2e}

Table~\ref{tab:e2e} applies inequality~\eqref{eq:breakeven-condition} to the
OpenAPI Specs dataset ($\Delta B = 6\,297$\,B without dedup;
$\Delta B = 32\,234$\,B with dedup) across representative network profiles.
CPU overhead: $\Delta t_{\text{enc}} = 0.038$\,ms,
$\Delta t_{\text{dec}} = 0.610$\,ms, total $= 0.648$\,ms.

\begin{table}[ht]
\centering
\caption{End-to-end latency estimate --- OpenAPI Specs (49\,KB).
         $t_{\text{total}} = t_{\text{enc}} + t_{\text{tx}} + t_{\text{dec}}$.
         Negative $\Delta t$ means SCON is faster end-to-end.}
\label{tab:e2e}
\begin{tabular}{l c rr rr r}
\toprule
Scenario & BW & \multicolumn{2}{c}{JSON (ms)} & \multicolumn{2}{c}{SCON min (ms)} & $\Delta t$ \\
         &    & enc+dec & tx & enc+dec & tx & (ms) \\
\midrule
\textbf{LAN}
  & 1\,Gbps   & 0.404 & 0.39 & 1.052 & 0.34 & $+0.60$ \\
\textbf{WAN}
  & 100\,Mbps & 0.404 & 3.93 & 1.052 & 3.42 & $+0.14$ \\
\textbf{Mobile}
  & 10\,Mbps  & 0.404 & 39.3 & 1.052 & 34.2 & $-4.39$ \\
\midrule
\multicolumn{7}{l}{\textit{With structural deduplication} (SCON dedup+min, 17\,KB):} \\
\midrule
LAN     & 1\,Gbps   & 0.404 & 0.39 & 1.052 & 0.13 & $+0.39$ \\
WAN     & 100\,Mbps & 0.404 & 3.93 & 1.052 & 1.35 & $-1.93$ \\
Mobile  & 10\,Mbps  & 0.404 & 39.3 & 1.052 & 13.5 & $-25.1$ \\
\bottomrule
\end{tabular}
\end{table}

\noindent \textbf{Key observations:}
\begin{enumerate}
\item Without deduplication, SCON breaks even below $\sim$100\,Mbps and
      provides a net gain on constrained links.
\item With structural deduplication, SCON is faster from WAN speeds
      downward ($-1.92$\,ms at 100\,Mbps).
\item On mobile links ($\leq$ 10\,Mbps), the transmission savings
      are an order of magnitude larger than the parsing penalty
      ($-25$\,ms with dedup).
\end{enumerate}

\noindent \textbf{Limitations of this model.}
Equation~\eqref{eq:breakeven} models only serialization cost and raw
transmission time. Real networks add round-trip latency (RTT), TLS
handshake, HTTP framing, and congestion-dependent throughput.
However, these factors are \emph{constant} across serialization
formats---both JSON and SCON traverse the same TCP/TLS/HTTP stack---so
they cancel in the differential $\Delta t$.
The model is therefore valid for comparing the \emph{marginal} effect
of format choice, not for predicting absolute request latency.

\subsection{Complementarity with Transport-Layer Compression}
\label{sec:gzip}

A common objection is that transport-layer compression (\texttt{gzip},
Brotli) already eliminates key redundancy in JSON, making structural
compaction redundant. While block compression does reduce the final
wire size, SCON's compaction operates at the \emph{semantic} level
(on the data tree) whereas \texttt{gzip} operates at the
\emph{syntactic} level (LZ77 sliding window over raw bytes).

This distinction has a measurable consequence: compression algorithms
scale their CPU and memory consumption with the size of the
uncompressed input buffer. By reducing the payload before it reaches
the transport layer, SCON reduces the computational work that
\texttt{gzip} must perform. Table~\ref{tab:gzip-cpu} quantifies
this effect.

\begin{table}[ht]
\centering
\caption{CPU time for \texttt{gzencode()} at level~6 (PHP~8, median
         of 500 iterations). Percentages indicate CPU savings vs
         compressing the JSON payload.}
\label{tab:gzip-cpu}
\begin{tabular}{l rrr rrr}
\toprule
Dataset
  & \multicolumn{3}{c}{gzip time (ms)}
  & \multicolumn{3}{c}{gzip output (B)} \\
  & JSON & SCON\textsubscript{min} & SCON\textsubscript{d+m}
  & JSON & SCON\textsubscript{min} & SCON\textsubscript{d+m} \\
\midrule
OpenAPI Specs
  & 0.171 & 0.153 \scriptsize($-$11\%)
  & 0.081 \scriptsize($-$53\%)
  & 1\,295 & 1\,227 & 1\,126 \\
Config Records
  & 0.459 & 0.422 \scriptsize($-$8\%)
  & 0.426 \scriptsize($-$7\%)
  & 4\,995 & 4\,832 & 4\,866 \\
DB Exports
  & 0.111 & 0.091 \scriptsize($-$18\%)
  & 0.087 \scriptsize($-$22\%)
  & 1\,473 & 1\,458 & 1\,458 \\
\bottomrule
\end{tabular}
\end{table}

\noindent SCON reduces gzip CPU time by 8--18\% without deduplication,
and up to 53\% with structural deduplication on schema-heavy payloads
(OpenAPI). In high-throughput API gateways where compression runs on
every response, these savings accumulate across thousands of requests
per second. SCON thus acts as a semantic pre-compressor: it reduces
the input that block compression must process, yielding compounding
benefits in both bandwidth and CPU budget.

SCON's smaller pre-compression footprint may also benefit constrained
environments (IoT, edge computing) where radio transmission time
(\emph{Time on Air}) and compression CPU budget are critical resources.
Quantifying this effect on specific radio protocols (LoRa, BLE) is
left as future work.

%% ============================================================
\section{Analysis of Implementation-Specific Overhead}
\label{sec:impl-overhead}

While Table~\ref{tab:rust-overhead} shows SCON decoding at 2.5--3.4$\times$ the
cost of \texttt{serde\_json}, a static audit of the Rust reference decoder
reveals that the dominant costs are avoidable memory allocations, not
grammar complexity. We identify three categories of \emph{accidental}
overhead in the current v0.x implementation.

\subsection{Profiled Hotspots}
\label{sec:hotspots}

\begin{table}[ht]
\centering
\caption{Decoder allocation hotspots identified by static code audit.
         Estimated impact ranges are first-order hypotheses, not measured.}
\label{tab:hotspots}
\begin{tabular}{l p{4.5cm} p{3.5cm} l}
\toprule
Location & Mechanism & Why removable & Status \\
\midrule
\texttt{decoder.rs:25}
  & Unconditional \texttt{input.to\_string()} clones entire input buffer
  & Non-minified input can be parsed as \texttt{\&str} directly
  & Patched \\
\texttt{decoder.rs:70}
  & Each parsed line stored as owned \texttt{String}
  & \texttt{Cow<'a, str>} enables zero-copy for unmodified lines
  & Not impl. \\
\texttt{decoder.rs:231}
  & Field names cloned per tabular row ($R \times K$ clones)
  & Shared \texttt{Arc<String>} or index-based mapping reduces to $K$ clones
  & Not impl. \\
\texttt{decoder.rs:622}
  & Delimiter split returns \texttt{Vec<String>} via \texttt{clone()}
  & \texttt{std::mem::take} or \texttt{Vec<\&str>} eliminates per-segment copies
  & Patched \\
\texttt{encoder.rs:355}
  & Character-by-character escape with per-char \texttt{push()}
  & Bulk \texttt{memcpy} for non-escaped runs
  & Not impl. \\
\bottomrule
\end{tabular}
\end{table}

\noindent For context, \texttt{serde\_json} has undergone over a decade of
optimization including SIMD-accelerated parsing and zero-copy deserialization.
The fact that a naive SCON implementation is within a small integer factor
(2.5--3.4$\times$) suggests that the format's computational cost is low
and potentially competitive with JSON at parity of implementation maturity.

\subsection{Preliminary Patch Results}
\label{sec:patch}

As an initial validation, we applied two patches to the Rust decoder:
(1)~eliminating the unconditional input clone, and
(2)~replacing \texttt{buffer.clone()} with \texttt{std::mem::take} in
delimiter splitting. Table~\ref{tab:patch} shows the decode ratio
before and after on the same hardware.

\begin{table}[ht]
\centering
\caption{Decode overhead ratio (SCON / serde\_json) before and after
         removing two allocation hotspots. 200 iterations, median.}
\label{tab:patch}
\begin{tabular}{l cc}
\toprule
Dataset & Before (v0.x) & After (patched) \\
\midrule
OpenAPI Specs   & 2.9$\times$ & 2.8$\times$ \\
Config Records  & 3.1$\times$$^\dagger$ & 2.5$\times$ \\
DB Exports      & 2.1$\times$ & 3.4$\times$$^\ddagger$ \\
\bottomrule
\end{tabular}

\smallskip
\noindent $^\dagger$ The pre-patch Config Records measurement was affected by
a decoder bug in expanded arrays with nested objects (only 1 of 40 items
decoded). The corrected decoder processes all items, so the before/after
comparison for this dataset is confounded with the bugfix.

\noindent $^\ddagger$ DB Exports showed high variance across runs due to
its small size (19\,KB); sub-millisecond timings are sensitive to
scheduling jitter.
\end{table}

\subsection{Internal Validity}
\label{sec:validity}

The estimated impacts in Table~\ref{tab:hotspots} are not additive:
optimizations may overlap (e.g., zero-copy lines subsume the input clone
benefit). Furthermore, static analysis cannot capture cache effects,
branch prediction, or allocator amortization. A comprehensive evaluation
requires CPU profiling (e.g., \texttt{perf stat}, flamegraphs) and
measurement of each optimization in isolation.

These projections will be validated in a follow-up revision
with per-optimization benchmarks.

%% ============================================================
\section{Optimization Roadmap}
\label{sec:roadmap}

Beyond the allocation hotspots addressed in Section~\ref{sec:impl-overhead},
several architectural optimizations remain unexplored:

\begin{description}
\item[Zero-copy parsing via \texttt{Cow<'a, str>}.]
  String values in SCON rarely require unescaping. A zero-copy decoder
  can return slices into the input buffer for unquoted strings and keys,
  avoiding allocation entirely for the common case.

\item[SIMD-accelerated indentation scanning.]
  Leading-space counting and semicolon scanning (minified mode) are
  amenable to SIMD vectorization (SSE4.2 / AVX2 / NEON), following
  the approach of \texttt{simdjson}~\cite{simdjson}.

\item[Canonical fast-path.]
  When the encoder emits with a known configuration (e.g., 1-space indent,
  no schemas), the decoder can activate a fast path that skips directive
  parsing, schema resolution, and flexible indent detection.

\item[Tabular batch decode.]
  Tabular arrays encode homogeneous rows. A specialized decoder can parse
  the header once and apply a fixed field layout to all subsequent rows,
  analogous to columnar batch processing.
\end{description}

\noindent These optimizations are orthogonal and composable. The Rust
implementation (Level~2) is the primary target for these improvements.

%% ============================================================
\section{Normative Semantics}
\label{sec:normative}

This section specifies behaviors that the EBNF grammar alone cannot
capture. Conforming implementations \textbf{MUST} follow these rules
(terminology per RFC~2119~\cite{rfc2119}).

\subsection{Duplicate Keys}

If an object contains two entries with the same key, the
\textbf{last occurrence wins}: the earlier value is silently replaced.
Implementations \textbf{SHOULD} emit a warning when duplicate keys
are detected during decoding. In \texttt{@@strict} mode, duplicate
keys \textbf{MUST} produce a parse error.

\subsection{Number Representation}

\begin{itemize}
\item Leading zeros are \textbf{not permitted} in integers
      (e.g., \texttt{007} is parsed as the string \texttt{"007"}).
\item \texttt{NaN}, \texttt{Infinity}, and \texttt{-Infinity} are
      \textbf{not valid} SCON values. Implementations that encounter
      these in source data \textbf{MUST} encode them as \texttt{null}
      or as quoted strings.
\item Floating-point values are serialized with sufficient decimal
      digits to guarantee round-trip fidelity (at least 15 significant
      digits for IEEE~754 double precision).
\end{itemize}

\subsection{Unicode and Encoding}

\begin{itemize}
\item SCON documents \textbf{MUST} be encoded in UTF-8 without BOM.
\item The \texttt{\textbackslash uXXXX} escape produces a single
      Unicode code point. Surrogate pairs (\texttt{U+D800}--\texttt{U+DFFF})
      \textbf{MUST NOT} appear as isolated escapes; implementations
      \textbf{SHOULD} reject them.
\item No Unicode normalization (NFC/NFKC) is applied by default.
      String equality is byte-level. Implementations \textbf{MAY}
      offer normalization as an option but \textbf{MUST NOT} apply
      it silently.
\end{itemize}

\subsection{Maximum Nesting Depth}

Implementations \textbf{MUST} support at least 64 levels of nesting.
Implementations \textbf{SHOULD} reject inputs exceeding a configurable
maximum depth (default: 512) to mitigate stack overflow and
billion-laughs-style attacks.

%% ============================================================
\section{Security Considerations}
\label{sec:security}

\subsection{Import Restrictions}

The \texttt{@use} directive loads external schema files. To prevent
Server-Side Request Forgery (SSRF) and file disclosure:

\begin{itemize}
\item Implementations \textbf{MUST NOT} resolve \texttt{http://},
      \texttt{https://}, or \texttt{file://} URIs by default.
      Only relative filesystem paths within a configurable base
      directory are permitted.
\item Implementations \textbf{MUST} reject import paths containing
      \texttt{..} path traversal sequences.
\item The maximum import depth \textbf{MUST} be bounded (default: 8).
      Circular imports \textbf{MUST} be detected and produce a
      parse error.
\end{itemize}

\subsection{Resource Limits}

To prevent denial-of-service via crafted inputs:

\begin{itemize}
\item \textbf{Maximum document size:} implementations \textbf{SHOULD}
      enforce a configurable limit (default: 64\,MB).
\item \textbf{Maximum nesting depth:} see Section~\ref{sec:normative}.
\item \textbf{Maximum schema count:} the \texttt{SchemaRegistry}
      \textbf{SHOULD} limit the number of registered schemas
      (default: 1\,024) to prevent memory exhaustion via
      \texttt{autoExtract} on adversarial inputs.
\item \textbf{Tabular row limit:} the declared row count in a
      tabular header \textbf{MUST} be validated against available
      data rows; implementations \textbf{MUST NOT} allocate memory
      based solely on the declared count.
\end{itemize}

\subsection{Directive Sandboxing}

The \texttt{@@} directive mechanism is extensible by design.
Implementations \textbf{MUST} ignore unrecognized directives
(forward compatibility). Directives \textbf{MUST NOT} execute
arbitrary code, modify filesystem state, or perform network
operations.

\bibliographystyle{plain}
\begin{thebibliography}{9}
\bibitem{rfc8259} T.~Bray, Ed.
  \textit{The JavaScript Object Notation (JSON) Data Interchange Format}.
  RFC~8259, IETF, December 2017.
\bibitem{iso14977} ISO/IEC~14977:1996.
  \textit{Information technology --- Syntactic metalanguage --- Extended BNF}.
\bibitem{yaml} O.~Ben-Kiki, C.~Evans, I.~d\"ot Net.
  \textit{YAML Ain't Markup Language (YAML) Version 1.2}.
  \url{https://yaml.org/spec/1.2/spec.html}, 2009.
\bibitem{scon-repo} G.~Moura.
  \textit{SCON: Structured Compact Object Notation --- Reference Implementations}.
  \url{https://github.com/QuijoteShin/scon}, 2026.
\bibitem{simdjson} G.~Langdale and D.~Lemire.
  \textit{Parsing Gigabytes of JSON per Second}.
  The VLDB Journal, 28(6):941--960, 2019.
\bibitem{rfc2119} S.~Bradner.
  \textit{Key words for use in RFCs to Indicate Requirement Levels}.
  RFC~2119, IETF, March 1997.
\end{thebibliography}

\end{document}
